\documentclass[10pt,a4paper]{amsart}
\usepackage{amsmath,amssymb,amsthm,graphicx}
% For commutative diagrams we use Xy-pic
\usepackage[all]{xy}

% We use 2cell for 2-commutative diagrams.
\xyoption{2cell}
% Theorem environments.
%
\theoremstyle{plain}
\newtheorem{theorem}[subsection]{Theorem}
\newtheorem{proposition}[subsection]{Proposition}
\newtheorem{lemma}[subsection]{Lemma}

\theoremstyle{definition}
\newtheorem{definition}[subsection]{Definition}
\newtheorem{example}[subsection]{Example}
\newtheorem{exercise}[subsection]{Exercise}
\newtheorem{situation}[subsection]{Situation}

\theoremstyle{remark}
\newtheorem{remark}[subsection]{Remark}
\newtheorem{remarks}[subsection]{Remarks}

\numberwithin{equation}{subsection}

% Macros
%
\def\lim{\mathop{\mathrm{lim}}\nolimits}
\def\colim{\mathop{\mathrm{colim}}\nolimits}
\def\Spec{\mathop{\mathrm{Spec}}}
\def\Hom{\mathop{\mathrm{Hom}}\nolimits}
\def\Ext{\mathop{\mathrm{Ext}}\nolimits}
\def\SheafHom{\mathop{\mathcal{H}\!\mathit{om}}\nolimits}
\def\SheafExt{\mathop{\mathcal{E}\!\mathit{xt}}\nolimits}
\def\Sch{\mathit{Sch}}
\def\Mor{\mathop{\mathrm{Mor}}\nolimits}
\def\Ob{\mathop{\mathrm{Ob}}\nolimits}
\def\Sh{\mathop{\mathit{Sh}}\nolimits}
\def\NL{\mathop{N\!L}\nolimits}
\def\CH{\mathop{\mathrm{CH}}\nolimits}
\def\proetale{{pro\text{-}\acute{e}tale}}
\def\etale{{\acute{e}tale}}
\def\QCoh{\mathit{QCoh}}
\def\Ker{\mathop{\mathrm{Ker}}}
\def\Im{\mathop{\mathrm{Im}}}
\def\Coker{\mathop{\mathrm{Coker}}}
\def\Coim{\mathop{\mathrm{Coim}}}

% Boxtimes
%
\DeclareMathSymbol{\boxtimes}{\mathbin}{AMSa}{"02}

%
% Macros for moduli stacks/spaces
%
\def\QCohstack{\mathcal{QC}\!\mathit{oh}}
\def\Cohstack{\mathcal{C}\!\mathit{oh}}
\def\Spacesstack{\mathcal{S}\!\mathit{paces}}
\def\Quotfunctor{\mathrm{Quot}}
\def\Hilbfunctor{\mathrm{Hilb}}
\def\Curvesstack{\mathcal{C}\!\mathit{urves}}
\def\Polarizedstack{\mathcal{P}\!\mathit{olarized}}
\def\Complexesstack{\mathcal{C}\!\mathit{omplexes}}
% \Pic is the operator that assigns to X its picard group, usage \Pic(X)
% \Picardstack_{X/B} denotes the Picard stack of X over B
% \Picardfunctor_{X/B} denotes the Picard functor of X over B
\def\Pic{\mathop{\mathrm{Pic}}\nolimits}
\def\Picardstack{\mathcal{P}\!\mathit{ic}}
\def\Picardfunctor{\mathrm{Pic}}
\def\Deformationcategory{\mathcal{D}\!\mathit{ef}}
\usepackage{cleveref}
\crefname{originallemma}{original lemma}{original lemmas}
\Crefname{originallemma}{Original lemma}{Original lemmas}
\makeatletter
\newlabel{algebra-lemma-pullback-module}{{02M1}{source}{}{sourcepointer-1}{}}
\newlabel{algebra-lemma-pullback-module@cref}{{[originallemma][0][]02M1}{[1][0][]source}{}{}{}}
\makeatother
\crefname{originallemma}{original lemma}{original lemmas}
\Crefname{originallemma}{Original lemma}{Original lemmas}
\makeatletter
\newlabel{lemma-additivity-divisors-restricted}{{02QG}{source}{}{sourcepointer-2}{}}
\newlabel{lemma-additivity-divisors-restricted@cref}{{[originallemma][0][]02QG}{[1][0][]source}{}{}{}}
\makeatother
\crefname{originallemma}{original lemma}{original lemmas}
\Crefname{originallemma}{Original lemma}{Original lemmas}
\makeatletter
\newlabel{lemma-flat-inverse-image-dimension}{{02R8}{source}{}{sourcepointer-3}{}}
\newlabel{lemma-flat-inverse-image-dimension@cref}{{[originallemma][0][]02R8}{[1][0][]source}{}{}{}}
\makeatother
\crefname{originallemma}{original lemma}{original lemmas}
\Crefname{originallemma}{Original lemma}{Original lemmas}
\makeatletter
\newlabel{situation-setup}{{02QL}{source}{}{sourcepointer-4}{}}
\newlabel{situation-setup@cref}{{[originallemma][0][]02QL}{[1][0][]source}{}{}{}}
\makeatother
\title{Flat pullback of the divisor of a rational function is the sum of the divisors on the irreducible components of the source}
\author{The Stacks Project Authors}
\date{stacks-project:chow.tex:local-snapshot-20261012; GNU FDL 1.2 or later}
\begin{document}
\maketitle
\noindent\textit{Editorial note.} All delivered mathematical text is the author's own, copied from the pinned Stacks Project chapter file chow.tex; no statement or proof step was rewritten, reordered or removed.\par\medskip
\noindent\textit{Editorial note.} Preamble provenance: the pinned chapter file opens with \textbackslash{}input\{preamble\} and carries no preamble text of its own, so the project's single Stacks-wide preamble staged at sources/stacks-chow/preamble.tex is excerpted verbatim as two spans, lines 14-18 (Xy-pic setup) and lines 154-217 (the native theorem declarations and mathematical macros). Nothing was added, removed or redefined; the wrapper class and the amsmath/amssymb/amsthm/graphicx base packages are supplied by the assembler.\par\medskip
\noindent\textit{Editorial note.} Same-chapter pointers are located at their \textbackslash{}label lines inside the pinned chapter file; cross-chapter pointers are located in the project's own tag list staged in the snapshot (sources/stacks-chow/tags.txt), at the verbatim lines carrying their canonical tags.\par\medskip
\noindent\textit{Allowed standard prerequisites (editorial).}\begin{itemize}
\item Dimension of the image of a cycle under a flat morphism. Same chapter chow.tex, label lemma-flat-inverse-image-dimension, invoked at proof line 3118. Gives dim\_delta(Z') at least k - 1 for the closure Z' of the image of an integral Z of dimension k + r - 1, which is what reduces the proof to the two cases Z' = Y and Z' a prime divisor.
\item Lengths and flat pullback of modules over a flat local homomorphism. Algebra, label algebra-lemma-pullback-module, invoked at proof line 3160. Identifies length\_A(A/(x)) times length\_B(B/m\_A B) with length\_B(B/xB), which is the form in which the desired equality is proved.
\item Additivity of divisors under a restricted family. Same chapter chow.tex, label lemma-additivity-divisors-restricted, invoked at proof line 3169. Supplies the final equality expressing length\_B(B/xB) as the sum over the minimal primes q\_i of the residue length times the local length; the reduction to that equality is carried out in the counted span.
\end{itemize}\medskip
\ifdefined\hypertarget\hypertarget{sourcepointer-1}{}\fi
\noindent\textit{Source reference.} Original lemma 02M1: sources/stacks-chow/tags.txt, source lines 3428--3428. Navigation pointer for \textbackslash{}ref\{algebra-lemma-pullback-module\} used inside the counted span (line 3160); the argument that consumes it is retained in the span.\par\medskip
\ifdefined\hypertarget\hypertarget{sourcepointer-2}{}\fi
\noindent\textit{Source reference.} Original lemma 02QG: sources/stacks-chow/pinned-original-source.tex, source lines 373--373. Navigation pointer for \textbackslash{}ref\{lemma-additivity-divisors-restricted\} used inside the counted span (line 3169); the argument that consumes it is retained in the span.\par\medskip
\ifdefined\hypertarget\hypertarget{sourcepointer-3}{}\fi
\noindent\textit{Source reference.} Original lemma 02R8: sources/stacks-chow/pinned-original-source.tex, source lines 2127--2127. Navigation pointer for \textbackslash{}ref\{lemma-flat-inverse-image-dimension\} used inside the counted span (line 3118); the argument that consumes it is retained in the span.\par\medskip
\ifdefined\hypertarget\hypertarget{sourcepointer-4}{}\fi
\noindent\textit{Source reference.} Original lemma 02QL: sources/stacks-chow/pinned-original-source.tex, source lines 1372--1372. Navigation pointer for \textbackslash{}ref\{situation-setup\} used inside the counted span (line 3099); the argument that consumes it is retained in the span.\par\medskip
\begin{lemma}
\label{lemma-prepare-flat-pullback-rational-equivalence}
Let $(S, \delta)$ be as in Situation \ref{situation-setup}.
Let $X$, $Y$ be schemes locally of finite type over $S$.
Assume $Y$ integral with $\dim_\delta(Y) = k$.
Let $f : X \to Y$ be a flat morphism of
relative dimension $r$. Then for $g \in R(Y)^*$ we have
$$
f^*\text{div}_Y(g) =
\sum n_j i_{j, *}\text{div}_{X_j}(g \circ f|_{X_j})
$$
as $(k + r - 1)$-cycles on $X$ where the sum is over the irreducible
components $X_j$ of $X$ and $n_j$ is the multiplicity of $X_j$ in $X$.
\end{lemma}

\begin{proof}
Let $Z \subset X$ be an integral closed subscheme of $\delta$-dimension
$k + r - 1$. We have to show that the coefficient $n$ of $[Z]$ in
$f^*\text{div}(g)$ is equal to the coefficient
$m$ of $[Z]$ in $\sum i_{j, *} \text{div}(g \circ f|_{X_j})$.
Let $Z'$ be the closure of $f(Z)$ which is an integral closed
subscheme of $Y$. By Lemma \ref{lemma-flat-inverse-image-dimension}
we have $\dim_\delta(Z') \geq k - 1$. Thus either $Z' = Y$
or $Z'$ is a prime divisor on $Y$. If $Z' = Y$, then the coefficients
$n$ and $m$ are both zero: this is clear for $n$ by definition
of $f^*$ and follows for $m$ because $g \circ f|_{X_j}$ is
a unit in any point of $X_j$ mapping to the generic point of $Y$.
Hence we may assume that $Z' \subset Y$ is a prime divisor.

\medskip\noindent
We are going to translate the equality of $n$ and $m$ into algebra.
Namely, let $\xi' \in Z'$ and $\xi \in Z$ be the generic points.
Set $A = \mathcal{O}_{Y, \xi'}$ and $B = \mathcal{O}_{X, \xi}$.
Note that $A$, $B$ are Noetherian, $A \to B$ is flat, local,
$A$ is a domain, and $\mathfrak m_AB$ is an ideal of definition
of the local ring $B$. The rational function $g$ is an element
of the fraction field $Q(A)$ of $A$.
By construction, the closed subschemes $X_j$
which meet $\xi$ correspond $1$-to-$1$ with minimal primes
$$
\mathfrak q_1, \ldots, \mathfrak q_s \subset B
$$
The integers $n_j$ are the corresponding lengths
$$
n_i = \text{length}_{B_{\mathfrak q_i}}(B_{\mathfrak q_i})
$$
The rational functions $g \circ f|_{X_j}$ correspond to the image
$g_i \in \kappa(\mathfrak q_i)^*$ of $g \in Q(A)$.
Putting everything together we see that
$$
n = \text{ord}_A(g) \text{length}_B(B/\mathfrak m_AB)
$$
and that
$$
m = \sum \text{ord}_{B/\mathfrak q_i}(g_i)
\text{length}_{B_{\mathfrak q_i}}(B_{\mathfrak q_i})
$$
Writing $g = x/y$ for some nonzero $x, y \in A$ we see that it suffices
to prove
$$
\text{length}_A(A/(x)) \text{length}_B(B/\mathfrak m_AB) =
\text{length}_B(B/xB)
$$
(equality uses Algebra, Lemma \ref{algebra-lemma-pullback-module})
equals
$$
\sum\nolimits_{i = 1, \ldots, s}
\text{length}_{B/\mathfrak q_i}(B/(x, \mathfrak q_i))
\text{length}_{B_{\mathfrak q_i}}(B_{\mathfrak q_i})
$$
and similarly for $y$. As $A \to B$ is flat it follows that $x$
is a nonzerodivisor in $B$. Hence the desired equality follows from
Lemma \ref{lemma-additivity-divisors-restricted}.
\end{proof}

\end{document}
